\documentclass{Group_Templ}

\begin{document}

\section*{Exercise 3.3}
\begin{proof}[s-t-shortest-path]
	Let $\delta(s,v)$ denote the shortest-path distance from $s$ to $v$.\newline
	\noindent\textbf{1. Upper Bound}\newline
	Consider any path $P = (s = v_0,\dots,v_k = t)$\newline
	For every edge along this path, the LP constraints imply:
	\begin{flalign*}
		d_{v1}&\le d_s+w(v_0,v_1) \\
		d_{v2}&\le d_{v_1}+w(v_1,v_2) \\
		&\vdots \\
		d_t&\le d_{v_{k-1}}+w(v_{k-1}, v_k)
	\end{flalign*}
	By summing these inequalities, the intermediate terms cancel. Since $d_s=0$, we obtain: \[d_t \le \sum_{i=1}^{k}w(v_{i-1},v_i)\]
	Since this holds for every path, including the shortest one:
	\[
	\boxed{
		d_t \le \delta(s,t)
	}\]
	\noindent\textbf{2. Feasibility}\newline
	We assign $d_v=\delta(s,v)$ for every vertex reachable from s.\newline
	For every edge $(u,v)$ between vertices reachable from $s$, the shortest-path property guarantees:
	\[
		\delta(s,v)\le \delta(s,u)+w(u,v)
	\]
	This holds because extending a shortest path to u by the edge $(u,v)$ cannot produce a shorter path than the shortest path to $v$.\newline
	
	\noindent For vertices unreachable from $s$, we choose a constant
	\[
	C = \max_{v \text{ reachable from } s} \delta(s,v)
	\]
	and assign $d_v=C$ to every unreachable vertex.
	
	\noindent There cannot be an edge from a reachable vertex to an
	unreachable vertex, since such an edge would make that
	vertex reachable.
	
	\noindent For edges between unreachable vertices, the constraints
	are satisfied because $C \leq C+w(u,v)$.
	
	\noindent For edges from unreachable to reachable vertices, the
	constraints are also satisfied since
	$\delta(s,v)\leq C \leq C+w(u,v)$.
	
	\noindent Therefore, all LP constraints are satisfied, and a feasible
	solution exists with
	\[
	d_t=\delta(s,t).
	\]
	\noindent\textbf{3. Conclusion}\newline
	Since $d_t$ cannot exceed $\delta(s,t)$, but this value can be achieved:
	\[
	\boxed{
		\max d_t = \delta(s,t)
	}\]
	Therefore, the LP correctly models the shortest-path
	problem, provided that $t$ is reachable from $s$.
\end{proof}
\newpage
\section*{Exercise 3.4}
Given:

\begin{flalign*}
	\min \quad x_1 - 2x_2 + x_3 + 4x_4,\\
	\text{s.t.}\quad
	x_1-x_2+2x_3+x_4 & = 20,\\
	x_1+2x_2+3x_3-x_4 &\geq -10,\\
	2x_1+4x_2-7x_3+x_4 &\leq 0,\\
	x_1 &\geq 4,\\
	x_2 &\geq 0,\\
	x_3 &\leq 0.
\end{flalign*}
The constraints $x_1-x_2+2x_3+x_4 = 20$, $2x_1+4x_2-7x_3+x_4 \leq 0$, $x_4 \in \mathbb{R}$, $x_1 \ge 4$ and $x_3 \leq 0$ have to be normalized.

\noindent\textbf{1. $x_1-x_2+2x_3+x_4 = 20$}:
\begin{flalign*}
	x_1-x_2+2x_3+x_4 &= 20 \Rightarrow x_1-x_2+2x_3+x_4 \ge 20 \text{ and } x_1-x_2+2x_3+x_4 \le 20 \\
	x_1-x_2+2x_3+x_4 &\le 20 \\
	\Leftrightarrow -x_1+x_2-2x_3-x_4 &\ge -20 \\
	\text{so } x_1-x_2+2x_3+x_4 &\ge 20 \text{ and } -x_1+x_2-2x_3-x_4 \ge -20
\end{flalign*}
\noindent\textbf{2. $2x_1+4x_2-7x_3+x_4 \leq 0$}:

\begin{flalign*}
	2x_1+4x_2-7x_3+x_4 &\leq 0 \\
	\Leftrightarrow -2x_1-4x_2+7x_3-x_4 &\geq 0
\end{flalign*}
\noindent\textbf{3. $x_4 \in \mathbb{R}$}:

\begin{flalign*}
x_4 &= x_4^+ - x_4^-,
\quad x_4^+, x_4^- \geq 0
\end{flalign*}
\noindent\textbf{4. $x_1 \ge 4$}:
Substitute:
\begin{flalign*}
x_1 &= 4 + y_1,\quad y_1 \ge 0
\end{flalign*}
\noindent\textbf{5. $x_3 \leq 0$}:
Substitute:
\begin{flalign*}
x_3 &= -y_3, \quad y_3 \ge 0
\end{flalign*}

\begin{flalign*}
	\min \quad (4+y_1) - 2x_2 - y_3 + 4(x_4^+ - x_4^-),\\
	\text{s.t.}\quad
	(4+y_1)-x_2-2y_3+(x_4^+ - x_4^-) &\ge 20,\\
	-(4+y_1)+x_2+2y_3-(x_4^+ - x_4^-) &\ge -20,\\
	(4+y_1)+2x_2-3y_3-(x_4^+ - x_4^-) &\geq -10,\\
	-(8+2y_1)-4x_2-7y_3-(x_4^+ - x_4^-) &\geq 0,\\
	y_1 &\geq 0,\\
	x_2 &\geq 0,\\
	y_3 &\ge 0, \\
	x_4^+ &\ge 0, \\
	x_4^- &\ge 0.
\end{flalign*}
\newpage
\noindent Expanding and rearranging the expressions yields:
\begin{flalign*}
	\min \quad y_1 - 2x_2 - y_3 + 4x_4^+ - 4x_4^-,\\
	\text{s.t.}\quad
	y_1-x_2-2y_3+x_4^+ - x_4^- &\ge 16,\\
	-y_1+x_2+2y_3-x_4^+ + x_4^- &\ge -16,\\
	y_1+2x_2-3y_3-x_4^+ + x_4^- &\geq -14,\\
	-2y_1-4x_2-7y_3-x_4^+ + x_4^- &\geq 8,\\
	y_1 &\geq 0,\\
	x_2 &\geq 0,\\
	y_3 &\ge 0, \\
	x_4^+ &\ge 0, \\
	x_4^- &\ge 0.
\end{flalign*}


\noindent\textbf{Reconstruction of the original solution:}

$z_{original}=z_{standard}+4$

\noindent Given a solution of the LP in standard form, the
original variables can be reconstructed as follows:
\begin{align*}
	x_1 &= 4+y_1,\\
	x_3 &= -y_3,\\
	x_4 &= x_4^+-x_4^-.
\end{align*}
The variable $x_2$ remains unchanged.

\section*{Exercise 3.5}

\begin{flalign*}
	x_1 = \text{Number of Strawberry Dream cups},\quad &x_2 = \text{Number of Chocolate Melt cups} \\
	\max \quad 5x_1+3x_2,\\
	\text{s.t.}\quad 2x_1+x_2&\le 50 \\
	x_1+3x_2&\le 30 \\
	x_1 \ge 0 \\
	x_2 \ge 0
\end{flalign*}

\end{document}
